Disaggregated systems require compact, queryable summaries of remote state to resolve the impedance mismatch between stateless compute and network-separated storage. We demonstrated that the system-level value of probabilistic data structures derives from two properties that resolve this mismatch. The first, dynamic mutability, enables state synchronization with evolving datasets through efficient element deletion. The second, algebraic composability, enables distributed aggregation without central coordination through operations that are commutative, associative, and idempotent.

The case studies on Cuckoo Filters (\S\ref{sec:membership}) and HyperLogLog (\S\ref{sec:cardinality}) validated this framework. We showed that specific deployment patterns emerge as necessary consequences of these properties: co-partitioning of filter state with data for oracles optimized for mutability, and hierarchical aggregation trees for oracles optimized for composability.

The framework (\S\ref{sec:framework}) establishes that mutability requires element-level information retention while composability requires lossy compression through commutative operations. This information-theoretic tension explains the absence of structures that exhibit both properties optimally. Structure selection follows from analyzing two workload characteristics: whether the dataset undergoes explicit deletion operations, and whether query results aggregate information from multiple partitions. The intersection of these characteristics determines which property the oracle must satisfy and thus which deployment pattern the system must instantiate.

Workloads requiring explicit element deletion and single-partition queries necessitate structures satisfying dynamic mutability, deployed in co-partitioned architectures where oracle state aligns with data partitioning. Workloads requiring cardinality aggregation across partitions over static or append-with-expiry datasets necessitate structures satisfying algebraic composability, deployed in hierarchical reduction architectures. Workloads requiring both properties must employ hybrid systems that maintain separate specialized oracles, as no single structure examined provides both properties optimally.

The framework exposes two research directions. The first is determining whether information-theoretic lower bounds prevent any structure from providing both constant-time deletion and idempotent merge operations without sacrificing space efficiency. The second is analyzing how security, privacy, and fault tolerance constraints interact with the core properties. Byzantine-resilient aggregation must filter malicious updates while preserving the composability required for distributed reduction. Differential privacy mechanisms must introduce noise in a manner that commutes with merge operations. Adaptive checkpointing protocols that bound error accumulation for composable oracles through threshold-based backup do not extend to structures lacking idempotent merge operations, creating an asymmetry in fault tolerance mechanisms.

The mutability-composability dichotomy explains why specific structures succeed in specific distributed contexts. Cuckoo Filters enable co-partitioned membership testing in systems where data undergoes deletion. HyperLogLog enables hierarchical cardinality aggregation in systems where dataset partitioning would otherwise require shuffling raw elements. The framework identifies which deployment constraints each property addresses and which architectural patterns each property enables. As systems continue to separate compute, memory, and storage into independent network-attached services, the framework provides a basis for analyzing which class of probabilistic oracle is appropriate for each layer of the disaggregated stack.
